% =============================================================================
% 004 / 068
% Status: DONE
% =============================================================================
% Notes:
%
% Source question:
% 你老能否闡釋一下環東海文明圈的概念，大概包括哪些地區？
% =============================================================================

\chapter[你老能否闡釋一下環東海文明圈的概念，大概包括哪些地區？]{你老能否闡釋一下環東海文明圈的概念，大概包括哪些地區？}
\qaanswer{
雖然我不喜歡用「自古以來」這個詞，但是受到「東中國沿岸寒流」「黑潮」、「臺灣暖流」、「對馬暖流」的季節性影響，東海自古以來就是一條有時限的海上高速公路。它連線了：山東半島、遼東半島、日本列島、朝鮮半島(高句麗、新羅、百濟)、沖繩(琉球)、臺灣(夷洲)、菲律賓(呂宋)。

這些地區的貿易據點構成的網路至少在4世紀初就成型了。貿易中涉及的物資、人員交往帶來了文化、文明的相互影響。尤其是受到東亞大陸政權影響相對更弱的日本、新羅、百濟、夷洲、呂宋與吳越的關係更為緊密。

從4世紀後半葉到10世紀末自有一套政治、經濟體系相對獨立的運作著並作用到以上各方。11世紀到14世紀由於宋帝國、元帝國先後對遠洋航海貿易的支援，這種關係得到了加強，先有南宋時期的宋日貿易，後有元帝國的「世界貿易」都是環東海文明圈的巔峰時期。之後由盛轉衰，進入到明帝國前期鄭和七下西洋後是一個長達近600年的漫長的海禁期，環東海文明戛然而止，少量文明交流的碎片被吳越武裝走私商、明政府軍、西葡荷商團僱傭軍、日本南部諸藩相互間的角力所裹挾，正常的大範圍的交流被明帝國的政策幹擾，直至清帝國時期全面中止。
}

